\documentclass[11pt]{article}
\usepackage[margin=1.1in]{geometry}
\usepackage{amsmath,amssymb,amsthm}
\usepackage{booktabs}
\usepackage[colorlinks=true,linkcolor=blue,citecolor=blue]{hyperref}

\newtheorem{lemma}{Lemma}
\newtheorem{proposition}{Proposition}
\newtheorem{corollary}{Corollary}
\theoremstyle{definition}
\newtheorem{definition}{Definition}
\theoremstyle{remark}
\newtheorem{remark}{Remark}

\newcommand{\HH}{\mathbb{H}}
\newcommand{\ZZ}{\mathbb{Z}}
\newcommand{\QQ}{\mathbb{Q}}
\newcommand{\CC}{\mathbb{C}}
\newcommand{\Res}{\operatorname{Res}}
\newcommand{\tk}{\widetilde{k}}
\newcommand{\KT}{\mathbf{K}_T}

\title{Eichler--Shimura defects of meromorphic period polynomials\\
on the winding-class lattice\\[4pt]
\large A computational session note}
\author{Claude (Fable~5), for D.~A.~McGady}
\date{July 3, 2026}

\begin{document}
\maketitle

\begin{abstract}
\noindent
This note documents one session of work on the canonical-contour problem (item T1 of the
pre-arXiv agenda): which winding classes of the two-segment contour give period polynomials
of meromorphic modular forms satisfying the Eichler--Shimura relations $r|(1+S)=0$ and
$r|(1+U+U^2)=0$.  Everything below concerns the period polynomial of the \emph{actual
meromorphic form}, computed by direct numerical integration of $f$ along the two-segment
contour; no pole subtraction or projection enters any computation.  The main outcomes:
(i)~two exact identities of the Hurwitz $T$-kernel --- a reflection identity forcing
$\KT|(1+S)=0$, and the fact that at integer $s$ the kernel is a Bernoulli polynomial ---
which together reduce the whole problem to integer linear algebra with closed-form entries;
(ii)~the ``the $|(1+S)$-defect of an $i$-pole is nonzero in every class'' conjecture is
\emph{false}: an explicit class kills it (verified to $10^{-33}$);
(iii)~for the weight-4 elliptic block, the classes satisfying \emph{both} relations are
completely characterized (Proposition~\ref{prop:joint}): a nonempty rank-4 lattice coset,
every weight in the defining equations derived by hand;
(iv)~in the clean class the original resonance proposition is restored and now
hand-derived: the pole breaks exactly its resonant relation, by exactly its residue.
Open items and all caveats are listed in \S\ref{sec:open}.
\end{abstract}

\section{Setting}\label{sec:conv}

Conventions follow the companion note~\cite{FCC}; this section makes the present note
self-contained.  Fix even weight $k$ and $n=k-2$, and let
$\Gamma=\mathrm{PSL}_2(\ZZ)=\langle S,T\rangle$ act on $\HH$ by $S\tau=-1/\tau$,
$T\tau=\tau+1$.  A \emph{meromorphic modular form} of weight $k$ is a quotient
$f=f_1/f_2$ with $f_i\in M_{k_i}$ holomorphic modular, $f_2\not\equiv0$, $k=k_1-k_2$; we
write $F_k$ for the space of these.  Let $\Gamma\!\cdot\!f^{-1}(\infty)$ denote the
$\Gamma$-orbit of the poles of $f$ in $\HH$.  For $\tau_0\in\HH$, the \emph{$S$- and
$T$-segments} are smooth paths in $\HH\setminus\Gamma\!\cdot\!f^{-1}(\infty)$,
\[
\gamma^S_{\tau_0}:\ -1/\tau_0\longrightarrow\tau_0,
\qquad
\gamma^T_{\tau_0}:\ \tau_0-1\longrightarrow\tau_0,
\]
and the \emph{$L$-integral} of $f\in F_k$ along the pair is
\begin{equation}\label{eq:Lint}
L^*(f,s)=e^{-\pi i s/2}\Big[\int_{\gamma^S_{\tau_0}} f(\tau)\,\tau^{s-1}\,d\tau
+\int_{\gamma^T_{\tau_0}} f(\tau)\,\tk(\tau,s)\,d\tau\Big],
\qquad
\tk(\tau,s)=\zeta(1-s,\tau+1)-e^{i\pi(s-1)}\,\zeta\!\big(1-(k-s),\tau+1\big),
\end{equation}
with $\zeta(s,x)$ the Hurwitz zeta function.  Both integrals are finite for every
$s\in\CC$ (compact contours avoiding the poles).

\begin{proposition}[Class dependence only~{\cite{FCC}}]\label{prop:indep}
$L^*(f,s)$ depends on $(\gamma^S_{\tau_0},\gamma^T_{\tau_0})$ only through the homotopy
class of the pair in $\HH\setminus\Gamma\!\cdot\!f^{-1}(\infty)$; within a class it is
independent of $\tau_0$ and of the paths.  For $f$ with no interior poles there is a
single class and $L^*$ is canonical.
\end{proposition}

\begin{proof}[Proof sketch]
Path independence within a class is Cauchy's theorem (the integrands are holomorphic off
the poles).  Independence of $\tau_0$ within a chamber follows because the boundary terms
of $\partial_{\tau_0}L^*$ cancel between the two segments, using only $T$-periodicity
$f(\tau_0-1)=f(\tau_0)$ and modularity $f(-1/\tau_0)=\tau_0^kf(\tau_0)$; see~\cite{FCC}.
\end{proof}

\begin{proposition}[Wall-crossing]\label{prop:wall}
Record a class by its \emph{winding data} $(X_S,X_T)$: integers, one pair per pole,
counting the net turns of $\gamma^S$, resp.\ $\gamma^T$, about each pole relative to a
reference class with value $L^*_0$.  Then
\begin{equation}\label{eq:wall}
L^*(f,s;X_S,X_T)=L^*_0(f,s)
+2\pi i\,e^{-\pi i s/2}\sum_{p}\Big[X_S(p)\,\Res_{\tau=p}\!\big(f\,\tau^{s-1}\big)
+X_T(p)\,\Res_{\tau=p}\!\big(f\,\tk(\tau,s)\big)\Big].
\end{equation}
\end{proposition}

\begin{proof}
Two pairs whose winding data differ by one unit at a single pole $p$ on one segment bound
a loop encircling $p$ once; the residue theorem adds $2\pi i$ times the residue of that
segment's integrand, and the prefactor $e^{-\pi is/2}$ multiplies through.  Sum over
crossings.
\end{proof}

The period polynomial (overall $(2\pi i)^{n+1}$ dropped, as in the
companion scripts) is
\begin{equation}\label{eq:rf}
r_f(X,Y)=\sum_{\ell=0}^{n}c_\ell\,X^{n-\ell}Y^\ell,
\qquad
c_\ell=(-1)^\ell\binom{n}{\ell}\,i^{\ell+1}\,L^*(f,\ell+1),
\end{equation}
equivalently $r_f=\int_{\gamma^S}f\,(X-\tau Y)^n\,d\tau+\int_{\gamma^T}f\,\KT\,d\tau$ with
the assembled kernel
\begin{equation}\label{eq:KT}
\KT(\tau;X,Y)=\sum_{\ell=0}^{n}\binom{n}{\ell}(-1)^{\ell}\,\tk(\tau,\ell+1)\,X^{n-\ell}Y^{\ell}.
\end{equation}
Slash actions: $(P|S)(X,Y)=P(-Y,X)$, $(P|U)(X,Y)=P(X-Y,X)$, $(P|U^2)(X,Y)=P(-Y,X-Y)$.
Define the two \emph{defect operators}
\[
D_S(P) := P\,|(1+S), \qquad D_U(P) := P\,|(1+U+U^2).
\]
For cusp forms and weakly holomorphic forms, $D_S(r_f)=D_U(r_f)=0$ (proven
in~\cite{FCC}; verified there against Brown's quasi-periods~\cite{Brown2017} and, via
Haberland's formula~\cite{Cohen2013}, against Petersson norms).
The question is what replaces this for meromorphic $f$.

\section{The winding lattice and the linear defect model}\label{sec:lattice}

By Propositions~\ref{prop:indep} and~\ref{prop:wall}, the value of $L^*$ depends on the
class only through the winding data $(X_S,X_T)$.  At the
integer values $s=\ell+1$ entering \eqref{eq:rf}, one crossing of pole $p$ by $\gamma^S$
shifts $r_f$ by $2\pi i\,R_p$, and one crossing by $\gamma^T$ shifts it by $2\pi i\,Q_p$,
where
\begin{equation}\label{eq:gens}
R_p:=\Res_{\tau=p}\!\big(f(\tau)\,(X-\tau Y)^n\big),
\qquad
Q_p:=\Res_{\tau=p}\!\big(f(\tau)\,\KT(\tau;X,Y)\big).
\end{equation}
Hence, writing $x=(x_p)$, $t=(t_p)$ for the winding integers of a class relative to a base
class with polynomial $r_f^{(0)}$,
\begin{equation}\label{eq:linear}
r_f(x,t) = r_f^{(0)} + 2\pi i\sum_p x_p\,R_p + 2\pi i\sum_p t_p\,Q_p ,
\end{equation}
and the defects $D_S,D_U$ are \emph{affine-linear in the winding integers}.  Every
question of the form ``is there a class with $D_S(r_f)=0$ and/or $D_U(r_f)=0$'' is
therefore an integer linear-algebra question with coefficient vectors
$D_S(R_p),D_U(R_p),D_S(Q_p),D_U(Q_p)$.  This specializes Proposition~\ref{prop:wall} to
integer $s$; the specialization was verified numerically to $32$ digits.

\section{Two exact identities of the $T$-kernel}\label{sec:identities}

\begin{lemma}[Reflection]\label{lem:reflect}
For even $k$ and integers $0\le m\le n$,
\[
\tk(\tau,\,k-1-m) \;=\; -(-1)^m\,\tk(\tau,\,m+1).
\]
\end{lemma}

\begin{proof}
At $s=m+1$ the phase is $e^{i\pi m}=(-1)^m$, so
$\tk(\tau,m+1)=\zeta(-m,\tau+1)-(-1)^m\zeta(m+2-k,\tau+1)$.
At $s=k-1-m$ the two Hurwitz arguments are $1-s=m+2-k$ and $1-(k-s)=-m$, and the phase is
$e^{i\pi(k-2-m)}=(-1)^m$ for even $k$; so
$\tk(\tau,k-1-m)=\zeta(m+2-k,\tau+1)-(-1)^m\zeta(-m,\tau+1)$,
which is $-(-1)^m$ times the previous display.
\end{proof}

\begin{corollary}\label{cor:KTS}
$\KT|S=-\KT$, hence $D_S(\KT)=0$ identically, and therefore:
\begin{enumerate}
\item[(i)] $D_S(Q_p)=0$ for every pole $p$ and every $f$: \emph{$T$-windings never move
the $|(1+S)$-defect}.
\item[(ii)] At $\tau_0=i$ the $S$-segment degenerates to a point, so
$r_f=\int_{\gamma^T}f\,\KT$ and $D_S(r_f)=\int_{\gamma^T}f\,D_S(\KT)=0$: the
$S$-relation holds automatically, for \emph{every} $f\in F_k$ with no pole on the
contour, in the $\tau_0=i$ class.
\end{enumerate}
\end{corollary}

\begin{proof}
In \eqref{eq:KT} substitute $(X,Y)\to(-Y,X)$ and re-index $\ell\to n-\ell$; the coefficient
of $X^{n-\ell}Y^{\ell}$ in $\KT|S$ is $\binom{n}{\ell}\tk(\tau,k-1-\ell)$, which by
Lemma~\ref{lem:reflect} equals $-(-1)^{\ell}\binom{n}{\ell}\tk(\tau,\ell+1)$, i.e.\ minus
the coefficient in $\KT$.  Items (i), (ii) follow because $D_S$ commutes with
$\Res_p(f\,\cdot)$ and with $\int_{\gamma^T}f\,\cdot$\,.
\end{proof}

Corollary~\ref{cor:KTS}(ii) upgrades an empirical observation (all five test forms held
$|(1+S)$ to $\sim10^{-33}$ at $\tau_0=i$, evidence item [E4] of the working ledger) to a
theorem, and (i) decouples the problem: the $S$-defect is controlled by $\gamma^S$-windings
alone, while $\gamma^T$-windings move only the $U$-defect.

\begin{lemma}[Bernoulli form of the kernel]\label{lem:bern}
At the integer points $s=m+1$, $0\le m\le n$, using $\zeta(-j,x)=-B_{j+1}(x)/(j+1)$,
\[
\tk(\tau,m+1) \;=\; -\frac{B_{m+1}(\tau+1)}{m+1} \;+\; (-1)^m\,\frac{B_{k-1-m}(\tau+1)}{k-1-m}\,,
\]
a polynomial in $\tau$.  In particular, for $k=4$:
\[
\tk(\tau,1)=-B_1(\tau+1)+\tfrac13B_3(\tau+1),\qquad
\tk(\tau,2)=-B_2(\tau+1),\qquad
\tk(\tau,3)=-\tk(\tau,1),
\]
\begin{equation}\label{eq:KT4}
\KT \;=\; \tk(\tau,1)\,(X^2-Y^2)\;+\;2\,B_2(\tau+1)\,XY .
\end{equation}
\end{lemma}

Consequently \emph{all} shift vectors \eqref{eq:gens} are residues of $f$ against explicit
polynomials: the canonical-class problem has closed-form integer data.  (The visible
$(X^2-Y^2)$ and $XY$ in \eqref{eq:KT4} are each killed by $(1+S)$ --- Corollary~\ref{cor:KTS}
made concrete.)

\section{The weight-4 elliptic block: hand derivations and verdicts}\label{sec:k4}

Let $f=E_4\Delta/E_6^2\in F_4$, the canonical block with a double pole at $i$ (and no
other poles mod $\Gamma$), Laurent data $a_{-2}=-1/(1728\pi^2)$ and the parity relation
of the following lemma.

\begin{lemma}[Parity at the $S$-fixed point]\label{lem:parity}
Let $f\in F_4$ have a double pole at $i$,
$f=a_{-2}(\tau-i)^{-2}+a_{-1}(\tau-i)^{-1}+O(1)$.  Then $a_{-1}=i\,a_{-2}$.
\end{lemma}

\begin{proof}
Write $\tau=i+\varepsilon$; then $-1/\tau=i-\varepsilon-i\varepsilon^2+O(\varepsilon^3)$,
so $f(-1/\tau)=a_{-2}\varepsilon^{-2}(1-2i\varepsilon)-a_{-1}\varepsilon^{-1}+O(1)$,
while $\tau^4f(\tau)=(1-4i\varepsilon)\big(a_{-2}\varepsilon^{-2}
+a_{-1}\varepsilon^{-1}\big)+O(1)$.  Modularity $f(-1/\tau)=\tau^4f(\tau)$ at order
$\varepsilon^{-1}$ gives $-2ia_{-2}-a_{-1}=a_{-1}-4ia_{-2}$, i.e.\ $a_{-1}=ia_{-2}$.
\end{proof}

\subsection{Residue polynomials at the $i$-orbit}
For a double pole, $\Res_i(fg)=a_{-1}g(i)+a_{-2}g'(i)$ for holomorphic $g$.  With
$g=(X-\tau Y)^2$:
\begin{equation}\label{eq:Ri}
R_i = a_{-1}(X-iY)^2 - 2a_{-2}Y(X-iY)
\;\overset{a_{-1}=ia_{-2}}{=}\;
i\,a_{-2}\,(X^2+Y^2).
\end{equation}
By $T$-periodicity of $f$, $R_{i\pm1}$ is \eqref{eq:Ri} with $X\to X\mp Y$:
\begin{equation}\label{eq:Ripm}
R_{i+1}=i\,a_{-2}\big((X-Y)^2+Y^2\big),\qquad
R_{i-1}=i\,a_{-2}\big((X+Y)^2+Y^2\big).
\end{equation}

\subsection{Defects of the shift vectors, by hand}
Direct substitution gives
\[
\begin{aligned}
(X^2+Y^2)\,\big|(1+S) &= 2\,(X^2+Y^2), &
\big((X\mp Y)^2+Y^2\big)\big|(1+S) &= 3\,(X^2+Y^2),\\[2pt]
(X^2+Y^2)\,\big|(1+U+U^2) &= 4\,(X^2-XY+Y^2), &
\big((X-Y)^2+Y^2\big)\big|(1+U+U^2) &= 4\,(X^2-XY+Y^2),\\
&&\big((X+Y)^2+Y^2\big)\big|(1+U+U^2) &= 8\,(X^2-XY+Y^2).
\end{aligned}
\]
Writing $\sigma:=2\pi i\,R_i=-2\pi a_{-2}(X^2+Y^2)$ for the unit on the $S$-side and
normalizing the $U$-side by $D_U(2\pi iR_i)$, the shift table is:

\begin{center}
\begin{tabular}{lcccccc}
\toprule
generator & $x_i$ & $x_{i-1}$ & $x_{i+1}$ & $t_i$ & $t_{i-1}$ & $t_{i+1}$\\
\midrule
$S$-defect weight (units of $\sigma$) & $2$ & $3$ & $3$ & $0$ & $0$ & $0$\\
$U$-defect weight (units of $D_U(2\pi iR_i)$) & $1$ & $2$ & $1$ & $7/12$ & $7/12$ & $19/12$\\
\bottomrule
\end{tabular}
\end{center}

\noindent
The zeros in the $T$-columns of the $S$-row are Corollary~\ref{cor:KTS}(i).  The integer
entries $2,3,3$ and $1,2,1$ are the one-line computations above.  The $T$-column entries
are derived in \S\ref{sec:Tcols}; every entry of the table is therefore exact.  (Each was
independently confirmed by PSLQ at 34 digits, off-ray residuals $\lesssim10^{-35}$.)

\subsection{The $T$-columns, by hand}\label{sec:Tcols}
From \eqref{eq:KT4}, $\Res_i(f\KT)=a_{-1}\KT(i)+a_{-2}\,\partial_\tau\KT(i)$, with
$\tk(\tau,1)=\tfrac13\tau^3+\tfrac12\tau^2-\tfrac56\tau-\tfrac12$ and
$B_2(\tau+1)=\tau^2+\tau+\tfrac16$, so that
$\tk(i,1)=-1-\tfrac76i$, $\tk'(i,1)=-\tfrac{11}6+i$, $B_2(i+1)=-\tfrac56+i$.
Using $a_{-1}=ia_{-2}$,
\[
Q_i \;=\; a_{-2}\Big[\big(i\,\tk(i,1)+\tk'(i,1)\big)(X^2-Y^2)
+ 2\big(iB_2(i+1)+(2i+1)\big)XY\Big]
\;=\; a_{-2}\Big[-\tfrac23\,(X^2-Y^2) \,+\, \tfrac73\,i\,XY\Big].
\]
Since $(X^2-Y^2)\,|(1+U+U^2)=0$ while $XY\,|(1+U+U^2)=X^2-XY+Y^2$,
\[
D_U(2\pi i\,Q_i) \;=\; -\tfrac{14}{3}\,\pi a_{-2}\,(X^2-XY+Y^2)
\;=\; \tfrac{7}{12}\;D_U(2\pi i\,R_i).
\]
For the translates, the Bernoulli difference equations give
$\KT(\tau+1)=\KT(\tau)+(\tau^2+2\tau)(X^2-Y^2)+4(\tau+1)XY$, whence (using
$1$-periodicity of $f$)
\[
Q_{i+1}-Q_i=\Res_i\!\Big(f\cdot\big[(\tau^2+2\tau)(X^2-Y^2)+4(\tau+1)XY\big]\Big)
= a_{-2}\big[\,i\,(X^2-Y^2)+4i\,XY\,\big],
\]
\[
Q_{i-1}-Q_i=-\Res_i\!\Big(f\cdot\big[(\tau^2-1)(X^2-Y^2)+4\tau\,XY\big]\Big)=0,
\]
the second vanishing by the same parity cancellation
($a_{-1}(-2)+a_{-2}(2i)=0$ and $a_{-1}(4i)+4a_{-2}=0$).  Hence
\[
\mu_U(t_{i-1})=\tfrac7{12},
\qquad
\mu_U(t_{i+1})=\tfrac7{12}+\frac{D_U\big(2\pi i(Q_{i+1}-Q_i)\big)}{D_U(2\pi i\,R_i)}
=\tfrac7{12}+1=\tfrac{19}{12},
\]
completing the table.  The equality $Q_{i-1}=Q_i$ --- the repeated $7/12$ --- is thus a
parity identity, not a numerical coincidence.

\subsection{The clean class: the resonance proposition, restored}
In the clean class ($\tau_0=0.3+1.4i$, straight segments, no pole enclosed), the measured
defects of the actual $r_f$ are
\begin{equation}\label{eq:clean}
r_f\,\big|(1+U+U^2)=0 \ \ (\text{residual }5.6\times10^{-48}),
\qquad
r_f\,\big|(1+S)= -2\pi a_{-2}\,(X^2+Y^2) \;=\; 2\pi i\,R_i ,
\end{equation}
the right-hand side now \emph{derived} in \eqref{eq:Ri} and equal to the measurement to
48 digits.  Together with the mirror behaviour of $\rho$-pole forms (which in their clean
classes hold $|(1+S)$ and break $|(1+U+U^2)$; see \S\ref{sec:rho}), this restores the
original proposition in class-anchored form: \emph{in the clean class, a pole on the
fixed locus of an elliptic element breaks exactly its resonant relation, by exactly its
residue polynomial.}  What is genuinely class-dependent --- and what the earlier
``in every class'' version missed --- is quantified next.

\subsection{The obstruction conjecture is false; joint classes exist}
The earlier conjecture (ledger item [E5]) asserted the $S$-defect of an $i$-pole is
nonzero in \emph{every} class, based on sweeping only $x_i$ (weight $2$: even shifts of an
odd base).  The table refutes it: $\gcd(2,3)=1$, and indeed with base defect $\sigma$,
\[
x=(1,-1,0):\quad \sigma\,(1 + 2 - 3) = 0,
\]
verified by direct high-precision recomputation of $r_f$ in that class:
$|(1+S)|/\mathrm{scale}=7\times10^{-33}$ (the $U$-relation then breaks, as the table
predicts).  More is true: \emph{both} relations can be killed simultaneously, and the
solution set can be characterized completely.

\begin{proposition}[Joint Eichler--Shimura exactness for the weight-4 block]\label{prop:joint}
Let $f=E_4\Delta/E_6^2$, and take as reference the clean class, whose base defect pair is
$\big(2\pi i R_i,\,0\big)$ \eqref{eq:clean}.  Among the depth-one winding classes
$(x_i,x_{i-1},x_{i+1};\,t_i,t_{i-1},t_{i+1})\in\ZZ^6$, those whose period polynomial
satisfies \emph{both} $r_f|(1+S)=0$ and $r_f|(1+U+U^2)=0$ are exactly the integer
solutions of
\begin{equation}\label{eq:jointsys}
1+2x_i+3x_{i-1}+3x_{i+1}=0,
\qquad
12x_i+24x_{i-1}+12x_{i+1}+7t_i+7t_{i-1}+19t_{i+1}=0 .
\end{equation}
This set is exactly the coset
\begin{equation}\label{eq:coset}
(x;t) \;=\; (1,-1,0;\,-1,0,1)
\;+\;\ZZ\Big\langle
(0,0,0;\,1,-1,0),\;
(0,0,0;\,-19,0,7),\;
(-3,2,0;\,1,0,-1),\;
(0,-1,1;\,-1,0,1)
\Big\rangle,
\end{equation}
a coset of a rank-$4$ sublattice of $\ZZ^6$.  In particular Eichler--Shimura exactness
holds in infinitely many classes, and exactness alone does not select a contour.
\end{proposition}

\begin{proof}
By the linear model \eqref{eq:linear} (Proposition~\ref{prop:wall} specialized to
integer $s$), the two defects of the class
$(x;t)$ are affine-linear in the winding integers with shift vectors
$D_S,D_U$ of $2\pi iR_p$, $2\pi iQ_p$ at $p\in\{i,i\pm1\}$.  All twelve weights are
derived above: the $S$-row is $(2,3,3;0,0,0)$ in units of $\sigma=2\pi iR_i$ (\S4.2 for the
integer entries via \eqref{eq:Ri}--\eqref{eq:Ripm}; the $T$-entries vanish by
Corollary~\ref{cor:KTS}(i)), and the $U$-row is
$(1,2,1;\tfrac7{12},\tfrac7{12},\tfrac{19}{12})$ in units of $D_U(2\pi iR_i)$ (\S4.2 and
\S\ref{sec:Tcols}).  Both defect components lie on the fixed rays $\CC\,(X^2{+}Y^2)$,
$\CC\,(X^2{-}XY{+}Y^2)$, so vanishing of the defect pair is equivalent to the two scalar
equations: base $1$ plus the $S$-row for the first; base $0$ plus the $U$-row, cleared of
denominators ($\times12$), for the second --- i.e.\ \eqref{eq:jointsys}.
Diophantine solution.  The integer solutions of the homogeneous equation $2a+3b+3c=0$
are exactly $\alpha(3,-2,0)+\gamma(0,-1,1)$: reducing mod $3$ forces $3\mid 2a$, hence
$3\mid a$, and the parametrization follows.  A particular solution of the first equation
of \eqref{eq:jointsys} is $x^{(0)}=(1,-1,0)$.  Writing $N:=x_i+2x_{i-1}+x_{i+1}$, the
second equation reads $7t_i+7t_{i-1}+19t_{i+1}=-12N$, with particular solution
$t=N\,(1,0,-1)$ (because $7-19=-12$) and homogeneous solutions $s(1,-1,0)+u(-19,0,7)$
(reducing mod $7$ forces $7\mid 19t_{i+1}$, hence $7\mid t_{i+1}$).  Since
$N(x^{(0)})=-1$, $N(-3,2,0)=1$, and $N(0,-1,1)=-1$, assembling particular-plus-kernel in
$x$ with the matching $t=N(1,0,-1)$ gives exactly the coset \eqref{eq:coset}; the four
displayed kernel vectors generate the \emph{full} kernel lattice, since any homogeneous
solution, after subtracting multiples of the third and fourth, has $x$-part zero and is
then a combination of the first two.  Each generator checks against
\eqref{eq:jointsys} directly, e.g.\ $(-3,2,0;1,0,-1)$: $-6+6=0$ and $-36+48+7-19=0$;
the particular solution: $1+2-3=0$ and $12-24-7+19=0$.
\end{proof}

\noindent
(A float box search found $66$ of these solutions with all windings in $[-3,3]$, each
consistent with \eqref{eq:jointsys} --- the numerical shadow of the proposition.  The
inputs a reader may wish to interrogate: the wall-crossing law, Proposition~\ref{prop:wall},
verified to $32$ digits; the parity Lemma~\ref{lem:parity}; and the clean-class base pair,
measured to $48$ digits, \eqref{eq:clean}.)

\begin{remark}[Why six winding integers, and not two]\label{rem:whysix}
One might expect just two encircling numbers --- one per segment --- for a form with a
single pole orbit.  But the puncture set of $\HH\setminus\Gamma\!\cdot\!f^{-1}(\infty)$ is
the \emph{whole orbit}, and windings about distinct orbit points are inequivalent because
the integrand weights are not translation-invariant: by \eqref{eq:Ripm},
$R_{i\pm1}\neq R_i$ --- the same pole of $f$, but a different residue polynomial, which is
where the distinct weights ($2$ vs.\ $3$; $7/12$ vs.\ $19/12$) come from.  The depth-one
window keeps the three orbit representatives adjacent to the contour, times two segments:
$\ZZ^6$.  The two-winding intuition is, in fact, precisely the truncation behind the
refuted conjecture of \S4.4: with windings at $i$ alone, the $S$-equation of
\eqref{eq:jointsys} collapses to $1+2x_i=0$, which has no integer solution --- the old
``parity obstruction'' --- and it is the translate windings, with their weight $3$, that
dissolve it.
\end{remark}

\section{The $\rho$-blocks and a generic pole}\label{sec:rho}

The same machinery was run on $E_6/j$ ($k=6$, triple pole at $\rho=e^{2\pi i/3}$),
$E_8/j=\Delta/E_4$ ($k=8$, simple pole at $\rho$), and the generic-pole form
$E_6/(j-j(1.1i))$ ($k=6$).  Base class $\tau_0=i$ throughout (which holds $|(1+S)$
identically by Corollary~\ref{cor:KTS}(ii)); windings allowed at the poles adjacent to the
contour (``depth one'': the poles and their $T$-translates and $S$-images next to the
fundamental domain), boxes $|x|\le3$.

\smallskip
\noindent\textbf{$E_6/j$.}  All $U$-defects lie on one real ray with PSLQ-exact weights
\[
\mu_U:\quad x_\rho:-11,\quad x_{\rho+1}:-3,\quad t_\rho:-\tfrac85,\quad t_{\rho+1}:-\tfrac{48}5,
\qquad\text{base}=1 .
\]
Killing $U$ alone is feasible: $3\cdot(-3)+1\cdot(-\tfrac85)+(-1)\cdot(-\tfrac{48}5)=-1$,
i.e.\ $(x,t)=\big((0,3),(1,-1)\big)$, verified to $4.1\times10^{-31}$ (the $S$-relation
then breaks, weight $6.0$).  Killing \emph{both} failed in the depth-one window: the two
$S$-defect vectors $D_S(R_\rho),D_S(R_{\rho+1})$ are $\ZZ$-independent (no rational ray),
pinning $x_\rho=x_{\rho+1}=0$; the $T$-windings alone then need
$\tfrac85t_1+\tfrac{48}5t_2=1$, i.e.\ $8t_1+48t_2=5$ --- impossible mod $8$.
\emph{Caveat: this is an obstruction of the truncated window, not (yet) of the full
winding lattice; deeper windings refine the reachable lattice and may dissolve it.}
No impossibility proposition is claimed here or anywhere in this note: the mod-$8$
argument is exact only within the depth-one window and rests on the PSLQ-identified
weights (not yet hand-derived, unlike the $k=4$ table).  The honest general statement is:

\smallskip
\noindent\textbf{Question.}
For which blocks does the winding lattice reach joint exactness --- equivalently, for
which pole orbits does the defect coset contain the origin?  The weight-$4$ $i$-block
does (Proposition~\ref{prop:joint}); the depth-one window of $E_6/j$ does not; the full
lattice for $E_6/j$, $\Delta/E_4$, and the generic pole is undecided.

\subsection{A second generic pole, at depth two: $z=0.47+0.99i$}\label{sec:z047}

To probe the generic case beyond depth one, take $f=E_6/(j-j(z))$ with $z=0.47+0.99i$
($|z|\approx1.096>1$, off the imaginary axis, off the edge ${\rm Im}\,\tau=1$: a
fundamental-domain point that avoids every special locus).  Because the base class
$\tau_0=i$ holds the $S$-relation identically (Corollary~\ref{cor:KTS}(ii)), one may
search using only moves that provably keep the $S$-defect at zero:
$T$-windings anywhere (Corollary~\ref{cor:KTS}(i)), and the pairs
``wind $q$ once, wind $Sq$ once backwards'', which are exactly $S$-silent since
$R_{Sq}=R_q|S$ for even $n$ (change of variables $\tau\to-1/\tau$ in \eqref{eq:gens},
using $f(-1/\tau)=\tau^kf(\tau)$ and $d(-1/\tau)=\tau^{-2}d\tau$).  The joint problem then reduces to a \emph{single} complex
linear Diophantine equation on the $U$-ray ($\dim_\CC\ker(1-U)=1$ at $n=4$):
$\sum_g x_g\mu_g=-1$ with $\mu_g:=D_U(\text{move}_g)/D_U(\text{base})$.

The computation (13 moves: 5 pairs, 8 $T$-windings, spanning orbit depth two; all checks
passed --- base $S$-defect $8\times10^{-45}$, pair $S$-defects $\lesssim10^{-40}$,
off-ray residuals $\lesssim10^{-41}$) shows what genericity does and does not destroy:

\begin{itemize}
\item The weights $\mu_g$ are now genuinely complex and (empirically) irrational ---
e.g.\ $\mu_{\mathrm{pair}(z,Sz)}=-26.9827\ldots+70.1532\ldots i$ --- as expected once
both the isotropy and the singular-moduli arithmetic of Remark~\ref{rem:iso} are absent.
\item The exact $\ZZ[\Gamma]$-module relations survive and are confirmed to working
precision (40 digits): PSLQ, asked for integer relations, first returns \emph{structural}
identities, e.g.
$\mu_{\mathrm{pair}(z,Sz)}=\mu_{T@Sz}-\mu_{T@(Sz-1)}$
and
$\mu_{T@z}+\mu_{T@Sz}=\mu_{T@(z-1)}+\mu_{T@(Sz-1)}$
--- Hurwitz difference-equation identities among transcendental numbers, holding on the
nose.  The module structure is real; only its \emph{arithmetic normalization} was special
at CM points.
\item No joint-kill class was found: the search (integer box on the 13 distinct weights;
then PSLQ with a greedy deflation of structural relations, maxcoeff $10^6$) terminated in
the relation
$8+3\mu_{T@Sz}+5\mu_{T@(z-1)}+2\mu_{T@(Sz-1)}+5\mu_{T@S(z+1)}=0$:
within this window the reachable integer multiples of the base $U$-defect include
$\mathbf{8}$, but no combination producing $1$ was found.
\end{itemize}

The recurrence of $8$ is striking: the depth-one obstruction of $E_6/j$ (\S5, mod-$8$)
and this depth-two generic-pole computation --- two entirely different pole locations at
the \emph{same weight} $k=6$ --- both stall at index $8$, while the weight-$4$ block was
jointly killable through its twelfths.  This raises the possibility that joint
reachability is governed by the \emph{weight} (through the Bernoulli arithmetic of the
kernel) rather than, or in addition to, the isotropy of the pole.  Caveats: the box was
narrow, the deflation path greedy, and a single PSLQ relation certifies reachability of
$8$ but not unreachability of $1$; ``index $8$'' is evidence, not a theorem.

\smallskip
\noindent\textbf{$\Delta/E_4$ ($k=8$).}  The $U$-defect space is genuinely
multi-dimensional (no single ray; $\dim\ker(1-U)=3$ on degree-$6$ forms), rank $7$ of
$10$ generators; no joint solution in the $|x|\le2$ window.  Inconclusive.

\smallskip
\noindent\textbf{$E_6/(j-j(1.1i))$ (generic pole).}  Base class: $S$-defect $=0$
(Corollary~\ref{cor:KTS}(ii)), $U$-defect $\ne0$.  Each relation is separately killable
(this reproduces the earlier routed-contour observations); rank $7$ of $12$ generators,
real solutions of the joint system exist, but no \emph{integer} joint solution appeared
with the five free windings in $[-3,3]$.  Inconclusive; see \S\ref{sec:open}.

\section{The corrected picture}\label{sec:picture}

For meromorphic $f$ the period polynomial is canonically a \emph{coset}: the set
\eqref{eq:linear} as $(x,t)$ ranges over the winding lattice.  The defect pair
$\big(D_S(r_f),D_U(r_f)\big)$ is an affine-linear image of that lattice; single defect
values are class data, the coset is the invariant.  The numerical record then organizes
itself:

\begin{itemize}
\item classes on the axis $D_S=0$ always exist --- e.g.\ $\tau_0=i$, now a theorem
(Cor.~\ref{cor:KTS}(ii));
\item classes on the axis $D_U=0$ existed in every case tested (clean class for the
$i$-block, \eqref{eq:clean}; routed classes for the $1.1i$-forms; the explicit solution
for $E_6/j$);
\item whether the coset contains the \emph{origin} --- both relations at once --- is the
only genuinely open obstruction.  For the weight-4 block it does
(Proposition~\ref{prop:joint}, a rank-4 lattice of such classes); for the $\rho$-blocks
and the generic pole it is undecided at depth one.
\end{itemize}

\begin{remark}[Isotropy versus singular moduli: what is actually special about $i$]\label{rem:iso}
Two independent features of the elliptic point do different work, and they should not be
bundled.  \emph{Isotropy}: $S$ fixes $i$, which (a) forces the parity relation
$a_{-1}=ia_{-2}$ between Laurent coefficients, (b) makes $X-iY$ an eigenvector of the
slash action ($(X-iY)|S=-i(X-iY)$), collapsing all defect vectors onto single rays, and
(c) is what turned Proposition~\ref{prop:joint} into a two-equation problem with
small-integer weights.  \emph{Singular moduli}: $i$ is a CM point, so the residue data is
arithmetically commensurable (algebraic multiples of powers of $\pi$) --- this is where
the clean rationals like $7/12$ come from.  At a generic pole both features die: no
parity constraint, generically a simple pole, and defect weights with no rational
structure (confirmed at $1.1i$ and at $0.47+0.99i$, \S\ref{sec:z047}).  What survives is
the exact $\ZZ[\Gamma]$-module structure --- for a simple pole the entire winding module
is generated by the \emph{single} polynomial $(X-zY)^n$ slashed around the group,
$R_{\gamma z}=R_z|\gamma^{-1}$, and its difference/orbit identities hold among
transcendental weights on the nose (\S\ref{sec:z047}).  The a-priori expectation was that
isotropy is the only mechanism that can create a genuine obstruction to joint exactness
(a torsion class in the stabilizer's cohomology), with generic poles unobstructed at
sufficient depth.  The weight-6 computations put this in tension: both the $\rho$-pole
and the generic pole stall at index $8$ in their windows, while the weight-4 $i$-block is
jointly killable --- suggesting the weight, through the Bernoulli arithmetic of the
kernel, may govern.  Unresolved.
\end{remark}

Two inequivalent notions of ``canonical class'' emerge, and they disagree whenever the
residue is nonzero:
\begin{enumerate}
\item \textbf{The clean class}: geometrically distinguished, essentially unique, defines
the $L$-integral for any meromorphic form; the Eichler--Shimura relations hold up to an
explicit anomaly which \emph{is} the residue polynomial \eqref{eq:clean} --- canonical
pole data in closed form.
\item \textbf{Eichler--Shimura--exact classes}: exact relations, but (where they exist) a
full lattice of them, each with a different $L^*$; exactness alone does not select a
contour, and per-block existence is an arithmetic question.
\end{enumerate}
Which notion the paper should adopt is a judgment call (the natural tie-breakers within
(2) being central-value crossing-invariance and the Haberland/Petersson pairing); the
data of this note supports stating (1) as the definition and (2) as structure theory of
the coset.

\section{Roadmap: nailing the canonical cohomology class}\label{sec:roadmap}

The following are the concrete steps that would turn the picture of \S\ref{sec:picture}
into a theorem selecting a canonical class (hence a canonical $L^*(f,s)$ for every
$f\in F_k$). They are ordered by dependency and by expected payoff.

\begin{enumerate}
\item[\textbf{R1.}] \textbf{Hand-derive the $E_6/j$ weight table} $(-11,-3,-\tfrac85,-\tfrac{48}5)$
from the Bernoulli form of the kernel (Lemma~\ref{lem:bern}) at $\rho$, exactly as the
$k=4$ table was derived in \S\ref{sec:Tcols}. This makes the mod-$8$ obstruction of
\S\ref{sec:rho} \emph{unconditional within its depth-one window} and removes the last
PSLQ-only input at $k=6$.
\item[\textbf{R2.}] \textbf{Run the depth-$\ge2$ test} (\texttt{t1\_depth2.py}, written but
never run) for the generic pole $E_6/(j-j(1.1i))$ and for $E_6/j$. This is the decisive
falsification of the free-module prediction (R4): if a joint-kill class appears at
depth two for the generic pole, the ``index-$8$'' stall is a truncation artifact and
generic poles are unobstructed; if it does not, the $H^2$ story is wrong.
\item[\textbf{R3.}] \textbf{Decide the index-$8$ pattern.} Is $8$ the true lattice index
for the depth-two window at $z=0.47+0.99i$ (\S\ref{sec:z047})? This needs the \emph{full}
relation lattice of the $\mu_g$, not a single PSLQ relation. Does the index drop with
depth? Do $k=8,10$ replace $8$ by something else --- i.e.\ is joint reachability governed
by the weight (Bernoulli arithmetic) rather than the pole?
\item[\textbf{R4.}] \textbf{Settle the $H^2$ obstruction.} Since
$\mathrm{PSL}_2(\ZZ)\cong\ZZ/2*\ZZ/3$, the obstruction to a jointly-exact class should
live in $H^2(\ZZ/2;M_f)\oplus H^2(\ZZ/3;M_f)$, $M_f$ the residue module. Two gaps to
close: (a) identify the two summands with the $S$- and $U$-defects; (b) pin the module
structure of the $T$-shift vectors $Q_p$ (Lemma~\ref{lem:bern} should make this explicit).
Prediction: free orbits (generic poles) $\Rightarrow$ no obstruction; elliptic-orbit
poles $\Rightarrow$ possible finite torsion. R2 tests this directly.
\item[\textbf{R5.}] \textbf{Map the $k=8$ ($\Delta/E_4$) geometry}, where the $U$-defect
space is genuinely multi-dimensional ($\dim\ker(1-U)=3$); the single-ray analysis does
not apply and the joint system must be treated as full-rank integer linear algebra.
\item[\textbf{R6.}] \textbf{Is the period polynomial unique on the coset?} For the
weight-$4$ block, do the (infinitely many) jointly-exact classes of
Proposition~\ref{prop:joint} give the \emph{same} $r_f$ modulo the space $W$ of rational
period polynomials, or genuinely different ones? If the former, ``the'' period polynomial
is well-defined even though the contour is not.
\item[\textbf{R7.}] \textbf{Make the canonical-class choice} (\S\ref{sec:picture}):
clean class vs.\ jointly-exact lattice. If the latter is chosen, a tie-breaker is needed
within the lattice; the candidates are central-value crossing-invariance and the
Haberland/Petersson pairing. This is the decision that upgrades ``a coset'' to ``a
canonical class.''
\end{enumerate}

\section{Open items and caveats}\label{sec:open}

\begin{enumerate}
\item \textbf{Hand derivation of the $E_6/j$ weight table} ($-11,-3,-8/5,-48/5$):
pending.  (The $k=4$ table is now fully derived, \S\ref{sec:Tcols}.)  This would make the
mod-$8$ window obstruction of \S\ref{sec:rho} unconditional within its window.
\item \textbf{Depth $\ge2$ windings}: run for the generic pole $0.47+0.99i$
(\S\ref{sec:z047}: no joint kill; index-$8$ relation found); still untested for
$E_6/(j-j(1.1i))$ and $E_6/j$ (\texttt{t1\_depth2.py}, never run).  All ``no joint
solution'' statements remain \emph{window-limited}, not proofs.
\item \textbf{The index-$8$ pattern at $k=6$}: is joint reachability governed by weight
(Bernoulli arithmetic of the kernel) rather than pole isotropy?  Sharp sub-questions: is
$8$ provably the lattice index for the depth-two window at $0.47+0.99i$ (requires the
full relation lattice, not one PSLQ relation); does the index drop with depth; and does
$k=8,10$ show $8\to$ something else.
\item \textbf{Speculative} (flagged as such; not verified): since
$\mathrm{PSL}_2(\ZZ)\cong\ZZ/2*\ZZ/3$, the obstruction to an Eichler--Shimura--exact class
should be a pair of torsion classes in $H^2(\ZZ/2;\,\cdot\,)\oplus H^2(\ZZ/3;\,\cdot\,)$
with coefficients in the winding module; free orbits (generic poles) would then carry no
obstruction, elliptic-orbit poles possibly a finite one.  Two identified gaps: the precise
identification of the two $H^2$ components with the two defects, and the module structure
of the $T$-shift vectors $Q_p$.  This predicts a depth-$\ge2$ joint solution for the
generic-pole form --- a falsifiable test, not yet run.
\item The ``$66$ solutions'' count is a float box search cross-checked against the exact
ray equations; the exact statement is the Smith-normal-form solution lattice of
\S\ref{sec:k4}.
\item The $k=8$ ($\Delta/E_4$) defect geometry (multi-dimensional $U$-space) is unmapped.
\end{enumerate}

\appendix
\section{Numerical methods and scripts}\label{app:methods}

All computations: \texttt{mpmath} at 34 significant digits, conventions imported from
\texttt{wip\_meromorphic\_periods.py}.  Period polynomials of the actual forms by direct
\texttt{mp.quad} of \eqref{eq:Lint} at $s=1,\dots,k-1$ (two-segment contour; $\tau_0=i$
degenerate class or $\tau_0=0.3+1.4i$ clean class).  Residue vectors \eqref{eq:gens} by
circle sums (radius $0.02$, $64$ points; Hurwitz zeta evaluated directly for the
$Q_p$).  Rational identification by PSLQ; integer feasibility by exact Smith normal form
on the cleared-denominator ray equations; joint searches by kernel-aware enumeration
cross-checked at full precision.  Session scripts (scratchpad, in order):
\texttt{t1\_lattice.py} (linear model, first sweeps), \texttt{t1\_structure.py} (exact ray
analysis, PSLQ weights, falsification of [E5], joint-kill box search),
\texttt{t1\_exact.py} (Smith-normal-form verdicts, extended $\rho$-reps, kernel-aware
search), \texttt{t1\_depth2.py} (depth-2 test, written, \emph{not run}),
\texttt{t1\_generic\_047.py} (the \S\ref{sec:z047} run: 40 digits, $N_T=90$ $q$-terms for
the low-${\rm Im}$ orbit points, exactly-$S$-silent move set, box + PSLQ-with-deflation
search).  Headline
residuals: $\KT|(1+S)$ machine-exact (identity); [E5]-falsifying class $7\times10^{-33}$;
clean-class $U$-relation $5.6\times10^{-48}$; $E_6/j$ $U$-kill $4.1\times10^{-31}$;
PSLQ off-ray residuals $\lesssim10^{-35}$.

\begin{thebibliography}{9}

\bibitem{FCC}
D.~A.~McGady,
\emph{[the finite-contour cocycles note]},
companion draft \texttt{finite\_contour\_cocycles\_short.tex}, in preparation (2026).

\bibitem{Brown2017}
F.~Brown, \emph{A class of non-holomorphic modular forms III},
arXiv:1710.07912 (2017).

\bibitem{Cohen2013}
H.~Cohen,
\emph{Haberland's formula and numerical computation of Petersson scalar products},
in \emph{ANTS X: Proceedings of the Tenth Algorithmic Number Theory Symposium},
Open Book Ser.\ \textbf{1}, Math.\ Sci.\ Publ., 2013, pp.\ 249--270.

\end{thebibliography}

\end{document}
